\section*{Capítulo 6}

\solucion{6.1} Posición larga, 1,000 barriles. Día 1: $(86.20-85.00)\times1{,}000=+1{,}200$. Día 2: $(84.70-86.20)\times1{,}000=-1{,}500$. Día 3: $(83.90-84.70)\times1{,}000=-800$. Total: $-1{,}100=(83.90-85.00)\times1{,}000$.

\solucion{6.2} $F=70\,e^{(0.04+0.01-0.08)\times0.5}=70\,e^{-0.015}=68.96$. Como $y>r+u$, el futuro está por debajo del spot: backwardation.

\solucion{6.3} $\Delta T=3/12=0.25$. $y=0.04-\ln(76/80)/0.25=0.04-(-0.05129)/0.25=0.04+0.2052=24.5\,\%$ anual.

\solucion{6.4} $h^*=0.85\times0.25/0.20=1.0625$. $N^*=1.0625\times600{,}000/1{,}000=637.5\approx638$ contratos. Fracción de varianza eliminada: $\rho^2=0.7225$, es decir, 72\,\%.

\solucion{6.5} En backwardation, todo precio a plazo está por debajo del spot. Al vender a plazo (swap a 89.38) se fija un precio menor que el actual (92.60): frente a hoy, $1.16\times(-3.22)=-3.7$ millones sobre el volumen del swap. Pero el inventario no puede venderse hoy; se venderá en uno a tres meses, y la curva indica que el mercado espera precios menores. Si el spot converge al futuro de enero, el inventario sin cobertura perdería 17.6 millones. Además, la cobertura elimina escenarios extremos (por ejemplo, $-80.6$ millones si el WTI cae 30\,\%) que una empresa sin liquidez no puede absorber.

\solucion{6.6} Ninguna. La cobertura sobre WTI paga en función del WTI, que no cambió. Las compras indexadas al Brent se encarecen 5.3 dólares por barril y ese costo adicional queda descubierto: es riesgo de base. Implica que el ratio de cobertura debe estimarse con la serie de costos de compra reales, no con el WTI, y que un ratio cercano a uno solo es válido si el diferencial es estable; si es volátil, la correlación cae y la efectividad de la cobertura también.
